\input{preamble.tex}
\usepackage{graphicx}

\begin{document}

\reversemarginpar

\pagestyle{fancy}
\fancyhead[L]{Terminale STMG}
\fancyhead[C]{\textbf{DS n°8a : variables aléatoires}}
\fancyhead[R]{\today}

\null\vspace{-30pt}
Nom / Prénom : \\

Consignes particulières : 
\begin{itemize}[label=$\bullet$]
	\item 
	La calculatrice est autorisée.
	\item 
	L'exercice \ref{exe:2a} peut être fait partiellement sur la feuille.
	\item 
	Toute trace de recherche est prise en compte.
\end{itemize}

\hrule

\exe{6}{
\includegraphics{code}
	On lance un dé équilibré à 6 faces. Si le résultat est 1 ou 2, on perd 8 euros. Si le résultat est 3, on gagne 2 euros. 
	Sinon, on gagne 5 euros. 
	Posons $X$ la variable aléatoire correspondant au gain du joueur.
	\begin{enumerate}
		\item
		Quelles sont les valeurs prises par $X$ ?
		\item
		Déterminer la loi de $X$.
		\item
		Calculer $E(X)$. Le jeu est-il équitable ? Justifier
	\end{enumerate}
}{exe:1a}{

}

\exe{8}{
Compléter le triangle de Pascal donné ci-dessous.

\begin{center}
\begin{tabular}{>{$}l<{$}|*{8}{c}}
\multicolumn{1}{l}{$n$} &&&&&&&&\\\cline{1-1} 
0 &\dots&&&&&&\\
1 &\dots&\dots&&&&&\\
2 &\dots&\dots&\dots&&&&\\
3 &\dots&\dots&\dots&\dots&&&\\
4 &\dots&\dots&\dots&\dots&\dots&&& \\ 
5 & \dots & \dots & \dots & \dots & \dots & \dots && \\
6 & \dots & \dots & \dots & \dots & \dots & \dots & \dots& \\ 
7 & \dots & \dots & \dots & \dots & \dots & \dots & \dots& \dots \\ \hline
\multicolumn{1}{l}{} &0&1&2&3&4&5&6&7\\\cline{2-9}
\multicolumn{1}{l}{} &\multicolumn{8}{c}{$k$}
\end{tabular}
\end{center}

Expliquer avec vos mots ou avec des formules pourquoi on observe des symétries dans les lignes de ce triangle.

}{exe:2a}{

}

\exe{6}{
On considère une variable aléatoire $X$ suivant une loi binomiale de paramètres $\left ( 5 ; 2/5 \right)$. Donner :
\begin{enumerate}
\item les valeurs possibles pour $X$ ;
\item la probabilité associée à chacune de ces valeurs ;
\item l'espérance de $X$.
\end{enumerate}
}{exe:3a}{

}

%%%%%%%%%%%%% SUJET B %%%%%%%%%%%%%%
\setcounter{Exercise}{0}

\newpage

\pagestyle{fancy}
\fancyhead[L]{Terminale STMG}
\fancyhead[C]{\textbf{DS n°8b : variables aléatoires}}
\fancyhead[R]{\today}

\null\vspace{-30pt}
Nom / Prénom : \\

Consignes particulières : 
\begin{itemize}[label=$\bullet$]
	\item 
	La calculatrice est autorisée.
	\item 
	L'exercice \ref{exe:2a} peut être fait partiellement sur la feuille.
	\item 
	Toute trace de recherche est prise en compte.
\end{itemize}

\hrule

\exe{6}{
	On lance un dé équilibré à 6 faces. Si le résultat est 1 ou 2, on perd 8 euros. Si le résultat est 3, on gagne 1 euros. 
	Sinon, on gagne 6 euros. 
	Posons $X$ la variable aléatoire correspondant au gain du joueur.
	\begin{enumerate}
		\item
		Quelles sont les valeurs prises par $X$ ?
		\item
		Déterminer la loi de $X$.
		\item
		Calculer $E(X)$. Le jeu est-il équitable ? Justifier
	\end{enumerate}
}{exe:1b}{

}

\exe{8}{
Compléter le triangle de Pascal donné ci-dessous.

\begin{center}
\begin{tabular}{>{$}l<{$}|*{8}{c}}
\multicolumn{1}{l}{$n$} &&&&&&&&\\\cline{1-1} 
0 &\dots&&&&&&\\
1 &\dots&\dots&&&&&\\
2 &\dots&\dots&\dots&&&&\\
3 &\dots&\dots&\dots&\dots&&&\\
4 &\dots&\dots&\dots&\dots&\dots&&& \\ 
5 & \dots & \dots & \dots & \dots & \dots & \dots && \\
6 & \dots & \dots & \dots & \dots & \dots & \dots & \dots& \\ 
7 & \dots & \dots & \dots & \dots & \dots & \dots & \dots& \dots \\ \hline
\multicolumn{1}{l}{} &0&1&2&3&4&5&6&7\\\cline{2-9}
\multicolumn{1}{l}{} &\multicolumn{8}{c}{$k$}
\end{tabular}
\end{center}

Expliquer avec vos mots ou avec des formules pourquoi on observe des symétries dans les lignes de ce triangle.

}{exe:2b}{

}

\exe{6}{
On considère une variable aléatoire $X$ suivant une loi binomiale de paramètres $\left ( 5 ; 2/3 \right)$. Donner :
\begin{enumerate}
\item les valeurs possibles pour $X$ ;
\item la probabilité associée à chacune de ces valeurs ;
\item l'espérance de $X$.
\end{enumerate}
}{exe:3b}{

}

%%%%%%%%%%%%%

%\newpage
%\fancyhead[C]{\textbf{Solutions}}
%\shipoutAnswer


\end{document}
